\chapter{Anexos}
\label{sec:anexos}

\section{Directivas y Parámetros de Simulación en LTspice}

Para la validación circuital en el software LTspice se emplearon los modelos SPICE provistos por los fabricantes para
el operacional TL081 (librería \texttt{TL081.301}) y modelos comparativos (\texttt{op07.lib}).

\begin{lstlisting}[style=ltspice, caption={Directivas SPICE empleadas para la caracterización en frecuencia y en el dominio temporal.}]
* Analisis de respuesta en frecuencia (Diagrama de Bode)
.ac dec 20 1m 100k

* Analisis transitorio con senal cardiaca sintetizada
.tran 0 2 0 100u

* Definicion de la fuente biopotencial ECG desde archivo de texto
V_ecg in_p in_m PWL file=ecg_source.txt

* Inclusion de modelos de amplificadores operacionales
.lib TL081.301
.lib op07.lib
\end{lstlisting}

\section{Protocolo de Calibración en Laboratorio}

Para la puesta a punto experimental del instrumento se recomienda el siguiente procedimiento secuencial:
\begin{enumerate}
  \item Verificación de polarización estática:
  Conectar las fuentes de alimentación partidas a $\pm \SI{12}{\volt}$ o $\pm \SI{15}{\volt}$. Cortocircuitar ambas
  entradas a masa ($e_1 = e_2 = \SI{0}{\volt}$) y comprobar con el multímetro que la tensión en la salida $V_o$ sea
  inferior a $\SI{10}{\milli\volt}$ en continua.
  \item Calibración de la ganancia de primera y segunda etapa:
  Inyectar en la entrada diferencial una señal senoidal de $\SI{10}{\milli\volt}$ de amplitud ($\SI{20}{\milli\volt\text{pp}}$) a $\SI{25}{\hertz}$. Con el
  osciloscopio monitoreando la salida de $U_3$, ajustar el trimpot $R_1$ hasta observar una ganancia intermedia de $44.67\ \text{V/V}$ ($\approx \SI{33.0}{\decibel}$, correspondiente a $\approx \SI{0.89}{\volt\text{pp}}$).
  \item Calibración de la ganancia global de salida:
  Trasladar la sonda del osciloscopio al terminal de salida final $V_o$. Sin alterar la señal de entrada, ajustar el
  trimpot de salida $R_9$ hasta obtener la ganancia total del sistema de diseño ($\approx 1336\ \text{V/V} \equiv \SI{62.5}{\decibel}$ a $\SI{25}{\hertz}$, con $V_o \approx \SI{26.7}{\volt\text{pp}}$ sin saturación dentro de la zona lineal de alimentación).
  \item Verificación con simulador cardíaco:
  Conectar el generador de bioseñales virtuales configurado en modo QRS adulto ($1\ \text{mVpp}$ a $60\ \text{bpm}$).
  Comprobar que en la pantalla del osciloscopio la señal exhiba una amplitud de $\SI{1}{\volt\text{pp}}$ con línea de base
  estable en cero voltios.
\end{enumerate}
